\subsection{The rational sawtooth in the M\"obius Gram form}
\label{subsec:endpoint-sawtooth-gram}

The arithmetic quantity in the classical Vasyunin Gram formula can
be expressed by an ordinarily convergent series at every rational
argument. We record the convention, the exact coefficient map, and
the part of the resulting signed sum that remains to be estimated.
Define the periodic function
\[
 B(x)=\begin{cases}1-2\{x\},&x\notin\mathbb Z,\\
                   0,&x\in\mathbb Z,
       \end{cases}
 \qquad
 g_B(h/k)=\sum_{\ell\ge1}\frac{B(\ell h/k)}{\ell}
 \quad(h,k\in\mathbb N).
\]
The value of \(g_B\) depends on the rational number, so its value
at \(dh/dk\) is its value at \(h/k\).

\begin{lemma}[Rational sawtooth and the Vasyunin sum]
\label{lem:endpoint-rational-sawtooth}
The defining series for \(g_B(h/k)\) converges in its displayed
order. For coprime positive integers \(h,k\), put
\[
 V(h,k)=\sum_{r=1}^{k-1}\left\{\frac{rh}{k}\right\}
                       \cot\frac{\pi r}{k}.
\]
Then
\begin{equation}\label{eq:endpoint-g-vasyunin}
 g_B(h/k)=-\frac{\pi}{k}V(h,k).
\end{equation}
The empty sum for \(k=1\) has value zero.
\end{lemma}

\begin{proof}
Reduce the fraction first. For a reduced denominator \(k\), the
sequence \(B(\ell h/k)\) is periodic of period \(k\), and its
sum over one period is
\[
 \sum_{r=1}^{k-1}(1-2r/k)=0.
\]
Its partial sums are bounded, so summation by parts proves convergence.
For \(k>1\), grouping the first \(Mk\) terms into residue classes
and using \(\psi(z+M)-\psi(z)=\sum_{j=0}^{M-1}(j+z)^{-1}\)
gives
\[
 \sum_{\ell=1}^{Mk}\frac{B(\ell h/k)}{\ell}
 =\frac1k\sum_{r=1}^{k-1}B(rh/k)
                   \bigl(\psi(M+r/k)-\psi(r/k)\bigr).
\]
Here \(\psi=\Gamma'/\Gamma\). Since
\(\psi(M+r/k)=\log M+O_k(M^{-1})\), the zero period sum gives
\[
 g_B(h/k)=-\frac1k\sum_{r=1}^{k-1}B(rh/k)\psi(r/k).
\]
Pair \(r\) with \(k-r\). The identities
\(B(-x)=-B(x)\) and
\(\psi(1-z)-\psi(z)=\pi\cot(\pi z)\) yield
\[
 g_B(h/k)=\frac{\pi}{2k}
             \sum_{r=1}^{k-1}B(rh/k)\cot(\pi r/k).
\]
The sum of \(\cot(\pi r/k)\) over \(1\le r<k\) is zero
by the same pairing. Substitution of
\(B(rh/k)=1-2\{rh/k\}\) proves
\eqref{eq:endpoint-g-vasyunin}. For \(k=1\), every summand in
the defining series vanishes.
\end{proof}

Let \(\vartheta_n(x)=\{1/(nx)\}\), and write
\[
 G_{hk}=\int_0^\infty\vartheta_h(x)\vartheta_k(x)\,dx,
 \qquad \kappa=\log(2\pi)-\gamma_E.
\]
The integral is finite: the integrand is bounded near zero and is
\(1/(hkx^2)\) for all sufficiently large \(x\).
For coprime \(h,k\), Vasyunin's formula is
\begin{equation}\label{eq:endpoint-vasyunin-coprime}
 G_{hk}=\frac\kappa2\left(\frac1h+\frac1k\right)
       +\frac{k-h}{2hk}\log\frac hk
       -\frac{\pi}{2hk}\bigl(V(h,k)+V(k,h)\bigr).
\end{equation}
We use this classical identity with the precise normalization in
B\'aez-Duarte--Balazard--Landreau--Saias
\cite[Proposition~89, p.~39; definition in Section~4, p.~16]
{bbls-autocorrelation-2003}. Their autocorrelation is
\(\mathcal A(\alpha)=\int_0^\infty\{u\}\{\alpha u\}u^{-2}du\);
the substitution \(u=1/(hx)\) gives
\(G_{hk}=h^{-1}\mathcal A(h/k)\), which converts their formula
exactly into \eqref{eq:endpoint-vasyunin-coprime}.
Their Proposition~88, p.~38, also gives the series form below,
with their Bernoulli series \(\varphi_1=-g_B/2\).

\begin{proposition}[The signed Gram contribution]
\label{prop:endpoint-signed-gram}
For all positive integers \(h,k\),
\begin{equation}\label{eq:endpoint-gram-sawtooth}
 G_{hk}=\frac\kappa2\left(\frac1h+\frac1k\right)
       +\frac{k-h}{2hk}\log\frac hk
       +\frac{g_B(h/k)}{2h}+\frac{g_B(k/h)}{2k}.
\end{equation}
For arbitrary real coefficients \(a_1,\ldots,a_N\), define
\[
 S_0=\sum_{n\le N}a_n,\quad S_1=\sum_{n\le N}\frac{a_n}{n},
 \quad L_0=\sum_{n\le N}a_n\log n,\quad
 L_1=\sum_{n\le N}\frac{a_n\log n}{n}.
\]
Then the full quadratic form is
\begin{equation}\label{eq:endpoint-full-signed-gram}
 a^{\mathsf T}Ga=S_0(\kappa S_1+L_1)-S_1L_0
           +\sum_{h,k\le N}\frac{a_ha_k}{h}g_B(h/k).
\end{equation}
In particular, for \(N\ge2\) and
\(a_n=\mu(n)w_N(n)\), \(w_N(n)=1-\log n/\log N\),
the last contribution is exactly
\begin{equation}\label{eq:endpoint-mobius-sawtooth}
 \sum_{h,k\le N}\frac{\mu(h)\mu(k)}h
             w_N(h)w_N(k)g_B(h/k).
\end{equation}
\end{proposition}

\begin{proof}
For coprime \(h,k\), substitute
\eqref{eq:endpoint-g-vasyunin} into
\eqref{eq:endpoint-vasyunin-coprime}. For general \(h=da\),
\(k=db\), where \((a,b)=1\), the change of variables \(y=dx\)
gives \(G_{da,db}=d^{-1}G_{ab}\). Every term on the right of
\eqref{eq:endpoint-gram-sawtooth} has the same scaling:
the ratios and the values of \(g_B\) are unchanged, and each
coefficient is multiplied by \(d^{-1}\). This proves
\eqref{eq:endpoint-gram-sawtooth} for every pair. The formula
\eqref{eq:endpoint-vasyunin-coprime} is used only for coprime pairs;
no convention for an unreduced-denominator \(V\) is needed.

In the finite quadratic form, interchanging \(h,k\) in the last
term of \eqref{eq:endpoint-gram-sawtooth} combines the two
\(g_B\)-terms into \(\sum a_ha_k g_B(h/k)/h\).
The first term sums to \(\kappa S_0S_1\). The logarithmic term is
\[
 \frac12\sum_{h,k\le N}a_ha_k
           (h^{-1}-k^{-1})(\log h-\log k)
 =S_0L_1-S_1L_0.
\]
These identities prove \eqref{eq:endpoint-full-signed-gram}.
Substitution of the stated real coefficients proves
\eqref{eq:endpoint-mobius-sawtooth}; all sums in this step are finite.
\end{proof}

The linear coefficient has an equally explicit sign convention.
For \(b_n=\langle\mathbf1_{(0,1)},\vartheta_n\rangle\),
\[
 b_n=\frac{1-\gamma_E+\log n}{n},\qquad
 \left\|\mathbf1_{(0,1)}-\sum_{n\le N}a_n\vartheta_n\right\|_2^2
       =1-2b^{\mathsf T}a+a^{\mathsf T}Ga.
\]
Indeed \(u=1/(nx)\) reduces the integral for \(b_n\) to
\(n^{-1}\int_{1/n}^\infty\{u\}u^{-2}du\).
The part below one is \(\log n\), while the integral from one
to an integer \(M\) is \(\log M-H_M+1\), tending to
\(1-\gamma_E\). Thus the convention
\(c_n=(\gamma_E-1-\log n)/n\) in
\cite[Section~17, p.~52]{caceres-margin-v2} is exactly \(c=-b\),
and its expression \(1+2c^{\mathsf T}a+a^{\mathsf T}Ga\)
is the displayed physical norm.

The long-numerator sum displayed by Maier--Rassias
\cite[Theorem~2.1, p.~3]{maier-rassias-2018} has indexing range
\(k^D\le h<2k^D\), \(D\ge2\), hence \(h\ge k^2\).
Its transpose has \(k\ge h^2\). The signed form
\eqref{eq:endpoint-mobius-sawtooth} also contains the balanced
region \(h<k^2\), \(k<h^2\). This range comparison uses only
the displayed indexing range; no quantitative estimate from that
source is invoked here. The comparable pairs occurring in the
compact M\"obius correlation satisfy
\(\lvert\log(a/b)\rvert\le2L\). When
\(\min(a,b)>e^{2L}\), they satisfy both \(a<b^2\) and
\(b<a^2\): for example, \(a\ge b\) gives
\(a\le e^{2L}b<b^2\), and \(b\le a<a^2\).
An estimate for this balanced signed arithmetic contribution, and
the required uniform bound on the analytically continued
\(R_\epsilon(f)\) near \(\epsilon=0\), remain open in the present
argument. The coprime series defining \(R_\epsilon\) is used in
its proved absolute-convergence region \(\Re\epsilon>1/2\);
its extension to \(\Re\epsilon>0\) retains the analytic-continuation
meaning already established above.
